% appendix_ablations.tex
% Supplementary ablation studies and variant sweeps moved from the main report.
% Include in Report_Latex_v2.tex with: % appendix_ablations.tex
% Supplementary ablation studies and variant sweeps moved from the main report.
% Include in Report_Latex_v2.tex with: % appendix_ablations.tex
% Supplementary ablation studies and variant sweeps moved from the main report.
% Include in Report_Latex_v2.tex with: % appendix_ablations.tex
% Supplementary ablation studies and variant sweeps moved from the main report.
% Include in Report_Latex_v2.tex with: \input{appendix_ablations}

\section{Ablation Studies and Variant Sweeps}
\label{app:ablations}

This appendix reports full sweep results for the ablation studies summarised in the main report.

% -------------------------------------------------------------------
\subsection{Phase 1: VAE KL-Weight Sweep}
\label{app:vae_sweep}

\begin{table}[H]
\centering
\small
\begin{tabularx}{\linewidth}{c X r r r}
\toprule
\textbf{KL weight} ($\beta$) & \textbf{Role} & \textbf{F1} & \textbf{AUROC} & \textbf{AUPRC} \\
\midrule
0.001 & Minimal regularisation  & \textbf{0.343} & 0.770 & 0.336 \\
0.005 & Balanced                & 0.340          & 0.772 & 0.372 \\
0.010 & Moderate regularisation & 0.335          & 0.766 & 0.369 \\
0.050 & Strong regularisation   & 0.302          & 0.750 & 0.329 \\
\bottomrule
\end{tabularx}
\caption{VAE KL-weight sweep. Increasing $\beta$ weakens reconstruction fidelity. All variants underperform the plain AE baseline (F1=0.468), confirming that the probabilistic objective does not help on this domain.}
\label{tab:vae_beta_sweep}
\end{table}

% -------------------------------------------------------------------
\subsection{Phase 1: Baseline AE Score Ablation}
\label{app:score_ablation}

\begin{figure}[H]
\centering
\includegraphics[width=0.95\linewidth]{images/autoencoder/ae64_score_ablation.png}
\caption{Score ablation on the baseline AE checkpoint: seven reduction rules applied to the same checkpoint produce F1 ranging from 0.24 to 0.47. Generated by \texttt{experiments/anomaly\_detection/autoencoder/x64/baseline/notebook.ipynb}.}
\label{fig:ae64_score_ablation}
\end{figure}

\begin{table}[H]
\centering
\small
\begin{tabularx}{\linewidth}{l X r}
\toprule
\textbf{Scoring Rule} & \textbf{Definition} & \textbf{F1} \\
\midrule
\texttt{topk\_abs\_mean} & Average top 20\% of pixel errors (sorted by magnitude) & \textbf{0.468} \\
\texttt{mse\_mean}       & Average squared pixel errors across entire image        & 0.410 \\
\texttt{max\_abs}        & Single worst (maximum absolute) pixel error             & 0.324 \\
\texttt{foreground\_mse} & Average error in foreground-masked region               & 0.317 \\
\texttt{mae\_mean}       & Average absolute pixel errors (MAE)                     & 0.300 \\
\texttt{pooled\_mae\_mean} & MAE with spatial pooling                              & 0.293 \\
\texttt{foreground\_mae} & MAE in foreground region                                & 0.239 \\
\bottomrule
\end{tabularx}
\caption{Score ablation: seven reduction rules on the same AE checkpoint. \texttt{topk\_abs\_mean} wins by $+0.057$ F1 over \texttt{mse\_mean}, confirming that selecting high-error pixels outperforms averaging all pixels.}
\label{tab:score_ablation}
\end{table}

\begin{figure}[H]
\centering
\includegraphics[width=0.95\linewidth]{images/autoencoder/score_ablation_summary.png}
\caption{Score ablation on x224 checkpoint: \texttt{topk\_abs\_mean} remains dominant at higher resolution, confirming that the pixel-selection insight holds across resolutions. Generated by \texttt{experiments/anomaly\_detection/autoencoder/x224/main/notebook.ipynb}.}
\label{fig:ae224_score_ablation}
\end{figure}

% -------------------------------------------------------------------
\subsection{Phase 3: Teacher-Student Backbone Comparison}
\label{app:ts_comparison}

\begin{table}[H]
\centering
\small
\caption{Teacher-student distillation results across all backbone variants.}
\label{tab:ts_results}
\begin{tabularx}{\linewidth}{>{\raggedright\arraybackslash}X c c c c}
\toprule
\textbf{Backbone} & \textbf{F1} & \textbf{AUROC} & \textbf{AUPRC} & \textbf{Notes} \\
\midrule
ResNet18 (topk20)              & 0.495          & 0.894          & 0.519          & Baseline \\
ResNet50 (mixed 2:1 topk20)    & \textbf{0.525} & 0.909          & 0.599          & Best CNN result \\
WideResNet50-2 layer2 (topk25) & 0.508          & 0.920          & 0.540          & Single-layer \\
WideResNet50-2 multilayer (topk15) & 0.524      & \textbf{0.923} & 0.546          & Multilayer \\
ViT-B/16 (x224)                & 0.163          & 0.661          & 0.106          & Architectural mismatch \\
\bottomrule
\end{tabularx}
\end{table}

\begin{figure}[H]
\centering
\includegraphics[width=\linewidth]{images/teacherstudent/ts_family_comparison_full.png}
\caption{Teacher-student distillation: all backbone variants compared across F1, AUROC, and AUPRC. WideResNet50-2 multilayer achieves the strongest AUROC (0.923); F1 is near-identical with ResNet50 (0.524 vs 0.525). Generated by \texttt{experiments/anomaly\_detection/report\_figures/notebook.ipynb}.}
\label{fig:ts_family_full}
\end{figure}

% -------------------------------------------------------------------
\subsection{Phase 3: FastFlow Architecture and Variants}
\label{app:fastflow}

\textbf{Stage 1: Feature extraction.} The input wafer map is converted from grayscale to RGB via channel repetition, resized to backbone input size via bilinear interpolation, then passed through a frozen WideResNet50-2 backbone. Features from layer2 (512 channels) and layer3 (2048 channels) are extracted and concatenated into a single tensor C.

\textbf{Stage 2: Normalising flow.} A flow with 4 sequential steps processes tensor C. Each step alternates between an orthogonal 1$\times$1 convolution (invertible channel mixing) and an affine coupling layer. The coupling layer splits channels in half, passes the first half through a small CNN (Conv3$\times$3 $\rightarrow$ GELU $\rightarrow$ Conv1$\times$1 $\rightarrow$ GELU $\rightarrow$ Conv3$\times$3, zero-initialised) to compute shift and log-scale parameters, then applies an affine transform to the second half.

\textbf{Stage 3: Anomaly scoring.} The transformed tensor Z is converted to a per-pixel anomaly map by computing $0.5 \times Z^2$ spatially. This heatmap is upsampled back to 64$\times$64 and reduced to a single wafer-level anomaly score via mean pooling.

\begin{figure}[H]
\centering
\includegraphics[width=\linewidth]{images/fastflow/fastflow_1.png}
\caption{FastFlow preprocessing and feature extraction pipeline. Generated by \texttt{experiments/anomaly\_detection/fastflow/x64/main/notebook.ipynb}.}
\label{fig:fastflow_pipeline}
\end{figure}

\begin{figure}[H]
\centering
\includegraphics[width=\linewidth]{images/fastflow/fastflow_2.png}
\caption{FastFlow normalising flow architecture: 4 sequential steps of orthogonal 1$\times$1 convolution and affine coupling. The final transformed tensor Z is converted to an anomaly map ($0.5 \times Z^2$). Generated by \texttt{experiments/anomaly\_detection/fastflow/x64/main/notebook.ipynb}.}
\label{fig:fastflow_flow}
\end{figure}

\begin{table}[H]
\centering
\small
\caption{FastFlow configuration variants on WideResNet50-2 backbone.}
\label{tab:fastflow_results}
\begin{tabularx}{\linewidth}{>{\raggedright\arraybackslash}X c c c c}
\toprule
\textbf{Configuration} & \textbf{F1} & \textbf{AUROC} & \textbf{AUPRC} & \textbf{Notes} \\
\midrule
Layer2+Layer3, 4 steps, mean      & \textbf{0.482} & \textbf{0.871} & \textbf{0.489} & Selected variant \\
Layer2+Layer3, 6 steps, mean      & 0.470          & 0.870          & 0.479          & More depth unhelpful \\
Layer2 only, 6 steps, topk-r0.15  & 0.442          & 0.884          & 0.460          & Single-layer weaker \\
\bottomrule
\end{tabularx}
\end{table}

\begin{figure}[H]
\centering
\includegraphics[width=\linewidth]{images/fastflow/variant_comparison_metrics.png}
\caption{FastFlow variant comparison: F1, AUROC, AUPRC across three configurations. Feature breadth (layer2+layer3) matters more than flow depth. Generated by \texttt{experiments/anomaly\_detection/fastflow/x64/main/notebook.ipynb}.}
\label{fig:fastflow_variants}
\end{figure}

% -------------------------------------------------------------------
\subsection{Phase 4: ViT-B/16 Block Depth Sweep}
\label{app:vit_block_sweep}

\begin{table}[H]
\centering
\caption{ViT-B/16 PatchCore block depth sweep at $224\!\times\!224$. All runs use identical
hyperparameters (5\% coreset, topk\_patch\_ratio 0.10, cosine distance). Block indices are
0-based over the 12-block encoder.}
\label{tab:vit_block_sweep}
\begin{tabular}{llcccc}
\toprule
Block & Depth    & F1             & AUROC          & AUPRC          & Recall \\
\midrule
3            & Early    & 0.550          & 0.934          & 0.622          & 0.732 \\
\textbf{6}   & \textbf{Mid} & \textbf{0.580} & \textbf{0.954} & 0.642     & \textbf{0.800} \\
9            & Mid-late & 0.580          & 0.943          & \textbf{0.677} & 0.788 \\
11           & Final    & 0.435          & 0.866          & 0.490          & 0.564 \\
\bottomrule
\end{tabular}
\end{table}

% -------------------------------------------------------------------
\subsection{Phase 4: DINOv2 ViT-B/14 Block Sweep}
\label{app:dinov2_block_sweep}

\begin{table}[H]
\centering
\caption{DINOv2 ViT-B/14 PatchCore block sweep at $224\!\times\!224$. All runs use identical hyperparameters.}
\label{tab:dinov2_block_sweep}
\begin{tabular}{llcccc}
\toprule
Block & Depth    & F1             & AUROC          & AUPRC          & Recall \\
\midrule
\textbf{6}   & \textbf{Mid} & \textbf{0.521} & \textbf{0.926} & 0.549     & \textbf{73.6\%} \\
9            & Mid-late & 0.492          & 0.915          & \textbf{0.561} & 68.4\% \\
11           & Final    & 0.443          & 0.885          & 0.504          & 58.0\% \\
\bottomrule
\end{tabular}
\end{table}

% -------------------------------------------------------------------
\subsection{Phase 4: MAE Fine-Tuning Setup and Per-Defect Results}
\label{app:mae_finetune}

\textbf{Setup.} The last 6 transformer blocks (blocks 6--11) and the final LayerNorm are fine-tuned for 10 epochs on the same 40k training normals used for the memory bank. A 2-layer MLP decoder (768 $\to$ 512 $\to$ patch pixels) produces the reconstruction target and is discarded after pre-training. AdamW with LR $=10^{-4}$ and cosine decay is used throughout. The 75\% masking run converges from loss 0.069 (epoch 1) to 0.042 (epoch 10); the 25\% masking run converges to 0.032 (epoch 10).

\begin{table}[H]
\centering
\caption{Frozen vs.\ MAE fine-tuned ViT-B/16 PatchCore at two masking ratios. All settings identical except backbone weights and masking ratio. All variants use the 95th-percentile validation-normal threshold.}
\label{tab:vit_mae_comparison}
\begin{tabular}{lcccc}
\toprule
Backbone & F1 & AUROC & AUPRC & Recall \\
\midrule
Frozen (ImageNet)          & 0.595          & 0.956          & 0.671          & ${\sim}$80.0\% \\
MAE fine-tuned (mask 75\%) & 0.595          & 0.959          & 0.717          & 82.4\% \\
MAE fine-tuned (mask 25\%) & \textbf{0.594} & \textbf{0.962} & \textbf{0.762} & \textbf{84.4\%} \\
\bottomrule
\end{tabular}
\end{table}

\begin{table}[H]
\centering
\caption{Per-defect recall for frozen ViT-B/16 vs.\ MAE fine-tuned variants.}
\label{tab:vit_mae_per_defect}
\begin{tabular}{lccc}
\toprule
\textbf{Defect type} & \textbf{Frozen} & \textbf{MAE 75\%} & \textbf{MAE 25\%} \\
\midrule
Scratch   & 0.727          & 0.455 & \textbf{0.636} \\
Loc       & \textbf{0.829} & 0.805 & \textbf{0.829} \\
Center    & 0.618          & 0.618 & \textbf{0.676} \\
Edge-Loc  & \textbf{0.705} & 0.682 & \textbf{0.705} \\
Edge-Ring & 0.941          & \textbf{0.971} & 0.961 \\
Donut     & 1.000          & 1.000 & 1.000 \\
Random    & 1.000          & 1.000 & 1.000 \\
Near-full & 1.000          & 1.000 & 1.000 \\
\midrule
\textbf{Overall recall} & ${\sim}$80.0\% & 82.4\% & \textbf{84.4\%} \\
\bottomrule
\end{tabular}
\end{table}

\textbf{Why 25\% masking outperforms 75\% on this domain.} The 75\% ratio is standard for natural images, which have dense, spatially redundant texture. Wafer maps are sparse binary grids with far less texture diversity. At 75\% masking, only 49 of 196 patches remain visible, providing insufficient local context for the model to learn meaningful wafer patch statistics. At 25\% masking, 147 patches remain visible per pass, giving the model richer local context and producing features better calibrated to wafer spatial structure.


\section{Ablation Studies and Variant Sweeps}
\label{app:ablations}

This appendix reports full sweep results for the ablation studies summarised in the main report.

% -------------------------------------------------------------------
\subsection{Phase 1: VAE KL-Weight Sweep}
\label{app:vae_sweep}

\begin{table}[H]
\centering
\small
\begin{tabularx}{\linewidth}{c X r r r}
\toprule
\textbf{KL weight} ($\beta$) & \textbf{Role} & \textbf{F1} & \textbf{AUROC} & \textbf{AUPRC} \\
\midrule
0.001 & Minimal regularisation  & \textbf{0.343} & 0.770 & 0.336 \\
0.005 & Balanced                & 0.340          & 0.772 & 0.372 \\
0.010 & Moderate regularisation & 0.335          & 0.766 & 0.369 \\
0.050 & Strong regularisation   & 0.302          & 0.750 & 0.329 \\
\bottomrule
\end{tabularx}
\caption{VAE KL-weight sweep. Increasing $\beta$ weakens reconstruction fidelity. All variants underperform the plain AE baseline (F1=0.468), confirming that the probabilistic objective does not help on this domain.}
\label{tab:vae_beta_sweep}
\end{table}

% -------------------------------------------------------------------
\subsection{Phase 1: Baseline AE Score Ablation}
\label{app:score_ablation}

\begin{figure}[H]
\centering
\includegraphics[width=0.95\linewidth]{images/autoencoder/ae64_score_ablation.png}
\caption{Score ablation on the baseline AE checkpoint: seven reduction rules applied to the same checkpoint produce F1 ranging from 0.24 to 0.47. Generated by \texttt{experiments/anomaly\_detection/autoencoder/x64/baseline/notebook.ipynb}.}
\label{fig:ae64_score_ablation}
\end{figure}

\begin{table}[H]
\centering
\small
\begin{tabularx}{\linewidth}{l X r}
\toprule
\textbf{Scoring Rule} & \textbf{Definition} & \textbf{F1} \\
\midrule
\texttt{topk\_abs\_mean} & Average top 20\% of pixel errors (sorted by magnitude) & \textbf{0.468} \\
\texttt{mse\_mean}       & Average squared pixel errors across entire image        & 0.410 \\
\texttt{max\_abs}        & Single worst (maximum absolute) pixel error             & 0.324 \\
\texttt{foreground\_mse} & Average error in foreground-masked region               & 0.317 \\
\texttt{mae\_mean}       & Average absolute pixel errors (MAE)                     & 0.300 \\
\texttt{pooled\_mae\_mean} & MAE with spatial pooling                              & 0.293 \\
\texttt{foreground\_mae} & MAE in foreground region                                & 0.239 \\
\bottomrule
\end{tabularx}
\caption{Score ablation: seven reduction rules on the same AE checkpoint. \texttt{topk\_abs\_mean} wins by $+0.057$ F1 over \texttt{mse\_mean}, confirming that selecting high-error pixels outperforms averaging all pixels.}
\label{tab:score_ablation}
\end{table}

\begin{figure}[H]
\centering
\includegraphics[width=0.95\linewidth]{images/autoencoder/score_ablation_summary.png}
\caption{Score ablation on x224 checkpoint: \texttt{topk\_abs\_mean} remains dominant at higher resolution, confirming that the pixel-selection insight holds across resolutions. Generated by \texttt{experiments/anomaly\_detection/autoencoder/x224/main/notebook.ipynb}.}
\label{fig:ae224_score_ablation}
\end{figure}

% -------------------------------------------------------------------
\subsection{Phase 3: Teacher-Student Backbone Comparison}
\label{app:ts_comparison}

\begin{table}[H]
\centering
\small
\caption{Teacher-student distillation results across all backbone variants.}
\label{tab:ts_results}
\begin{tabularx}{\linewidth}{>{\raggedright\arraybackslash}X c c c c}
\toprule
\textbf{Backbone} & \textbf{F1} & \textbf{AUROC} & \textbf{AUPRC} & \textbf{Notes} \\
\midrule
ResNet18 (topk20)              & 0.495          & 0.894          & 0.519          & Baseline \\
ResNet50 (mixed 2:1 topk20)    & \textbf{0.525} & 0.909          & 0.599          & Best CNN result \\
WideResNet50-2 layer2 (topk25) & 0.508          & 0.920          & 0.540          & Single-layer \\
WideResNet50-2 multilayer (topk15) & 0.524      & \textbf{0.923} & 0.546          & Multilayer \\
ViT-B/16 (x224)                & 0.163          & 0.661          & 0.106          & Architectural mismatch \\
\bottomrule
\end{tabularx}
\end{table}

\begin{figure}[H]
\centering
\includegraphics[width=\linewidth]{images/teacherstudent/ts_family_comparison_full.png}
\caption{Teacher-student distillation: all backbone variants compared across F1, AUROC, and AUPRC. WideResNet50-2 multilayer achieves the strongest AUROC (0.923); F1 is near-identical with ResNet50 (0.524 vs 0.525). Generated by \texttt{experiments/anomaly\_detection/report\_figures/notebook.ipynb}.}
\label{fig:ts_family_full}
\end{figure}

% -------------------------------------------------------------------
\subsection{Phase 3: FastFlow Architecture and Variants}
\label{app:fastflow}

\textbf{Stage 1: Feature extraction.} The input wafer map is converted from grayscale to RGB via channel repetition, resized to backbone input size via bilinear interpolation, then passed through a frozen WideResNet50-2 backbone. Features from layer2 (512 channels) and layer3 (2048 channels) are extracted and concatenated into a single tensor C.

\textbf{Stage 2: Normalising flow.} A flow with 4 sequential steps processes tensor C. Each step alternates between an orthogonal 1$\times$1 convolution (invertible channel mixing) and an affine coupling layer. The coupling layer splits channels in half, passes the first half through a small CNN (Conv3$\times$3 $\rightarrow$ GELU $\rightarrow$ Conv1$\times$1 $\rightarrow$ GELU $\rightarrow$ Conv3$\times$3, zero-initialised) to compute shift and log-scale parameters, then applies an affine transform to the second half.

\textbf{Stage 3: Anomaly scoring.} The transformed tensor Z is converted to a per-pixel anomaly map by computing $0.5 \times Z^2$ spatially. This heatmap is upsampled back to 64$\times$64 and reduced to a single wafer-level anomaly score via mean pooling.

\begin{figure}[H]
\centering
\includegraphics[width=\linewidth]{images/fastflow/fastflow_1.png}
\caption{FastFlow preprocessing and feature extraction pipeline. Generated by \texttt{experiments/anomaly\_detection/fastflow/x64/main/notebook.ipynb}.}
\label{fig:fastflow_pipeline}
\end{figure}

\begin{figure}[H]
\centering
\includegraphics[width=\linewidth]{images/fastflow/fastflow_2.png}
\caption{FastFlow normalising flow architecture: 4 sequential steps of orthogonal 1$\times$1 convolution and affine coupling. The final transformed tensor Z is converted to an anomaly map ($0.5 \times Z^2$). Generated by \texttt{experiments/anomaly\_detection/fastflow/x64/main/notebook.ipynb}.}
\label{fig:fastflow_flow}
\end{figure}

\begin{table}[H]
\centering
\small
\caption{FastFlow configuration variants on WideResNet50-2 backbone.}
\label{tab:fastflow_results}
\begin{tabularx}{\linewidth}{>{\raggedright\arraybackslash}X c c c c}
\toprule
\textbf{Configuration} & \textbf{F1} & \textbf{AUROC} & \textbf{AUPRC} & \textbf{Notes} \\
\midrule
Layer2+Layer3, 4 steps, mean      & \textbf{0.482} & \textbf{0.871} & \textbf{0.489} & Selected variant \\
Layer2+Layer3, 6 steps, mean      & 0.470          & 0.870          & 0.479          & More depth unhelpful \\
Layer2 only, 6 steps, topk-r0.15  & 0.442          & 0.884          & 0.460          & Single-layer weaker \\
\bottomrule
\end{tabularx}
\end{table}

\begin{figure}[H]
\centering
\includegraphics[width=\linewidth]{images/fastflow/variant_comparison_metrics.png}
\caption{FastFlow variant comparison: F1, AUROC, AUPRC across three configurations. Feature breadth (layer2+layer3) matters more than flow depth. Generated by \texttt{experiments/anomaly\_detection/fastflow/x64/main/notebook.ipynb}.}
\label{fig:fastflow_variants}
\end{figure}

% -------------------------------------------------------------------
\subsection{Phase 4: ViT-B/16 Block Depth Sweep}
\label{app:vit_block_sweep}

\begin{table}[H]
\centering
\caption{ViT-B/16 PatchCore block depth sweep at $224\!\times\!224$. All runs use identical
hyperparameters (5\% coreset, topk\_patch\_ratio 0.10, cosine distance). Block indices are
0-based over the 12-block encoder.}
\label{tab:vit_block_sweep}
\begin{tabular}{llcccc}
\toprule
Block & Depth    & F1             & AUROC          & AUPRC          & Recall \\
\midrule
3            & Early    & 0.550          & 0.934          & 0.622          & 0.732 \\
\textbf{6}   & \textbf{Mid} & \textbf{0.580} & \textbf{0.954} & 0.642     & \textbf{0.800} \\
9            & Mid-late & 0.580          & 0.943          & \textbf{0.677} & 0.788 \\
11           & Final    & 0.435          & 0.866          & 0.490          & 0.564 \\
\bottomrule
\end{tabular}
\end{table}

% -------------------------------------------------------------------
\subsection{Phase 4: DINOv2 ViT-B/14 Block Sweep}
\label{app:dinov2_block_sweep}

\begin{table}[H]
\centering
\caption{DINOv2 ViT-B/14 PatchCore block sweep at $224\!\times\!224$. All runs use identical hyperparameters.}
\label{tab:dinov2_block_sweep}
\begin{tabular}{llcccc}
\toprule
Block & Depth    & F1             & AUROC          & AUPRC          & Recall \\
\midrule
\textbf{6}   & \textbf{Mid} & \textbf{0.521} & \textbf{0.926} & 0.549     & \textbf{73.6\%} \\
9            & Mid-late & 0.492          & 0.915          & \textbf{0.561} & 68.4\% \\
11           & Final    & 0.443          & 0.885          & 0.504          & 58.0\% \\
\bottomrule
\end{tabular}
\end{table}

% -------------------------------------------------------------------
\subsection{Phase 4: MAE Fine-Tuning Setup and Per-Defect Results}
\label{app:mae_finetune}

\textbf{Setup.} The last 6 transformer blocks (blocks 6--11) and the final LayerNorm are fine-tuned for 10 epochs on the same 40k training normals used for the memory bank. A 2-layer MLP decoder (768 $\to$ 512 $\to$ patch pixels) produces the reconstruction target and is discarded after pre-training. AdamW with LR $=10^{-4}$ and cosine decay is used throughout. The 75\% masking run converges from loss 0.069 (epoch 1) to 0.042 (epoch 10); the 25\% masking run converges to 0.032 (epoch 10).

\begin{table}[H]
\centering
\caption{Frozen vs.\ MAE fine-tuned ViT-B/16 PatchCore at two masking ratios. All settings identical except backbone weights and masking ratio. All variants use the 95th-percentile validation-normal threshold.}
\label{tab:vit_mae_comparison}
\begin{tabular}{lcccc}
\toprule
Backbone & F1 & AUROC & AUPRC & Recall \\
\midrule
Frozen (ImageNet)          & 0.595          & 0.956          & 0.671          & ${\sim}$80.0\% \\
MAE fine-tuned (mask 75\%) & 0.595          & 0.959          & 0.717          & 82.4\% \\
MAE fine-tuned (mask 25\%) & \textbf{0.594} & \textbf{0.962} & \textbf{0.762} & \textbf{84.4\%} \\
\bottomrule
\end{tabular}
\end{table}

\begin{table}[H]
\centering
\caption{Per-defect recall for frozen ViT-B/16 vs.\ MAE fine-tuned variants.}
\label{tab:vit_mae_per_defect}
\begin{tabular}{lccc}
\toprule
\textbf{Defect type} & \textbf{Frozen} & \textbf{MAE 75\%} & \textbf{MAE 25\%} \\
\midrule
Scratch   & 0.727          & 0.455 & \textbf{0.636} \\
Loc       & \textbf{0.829} & 0.805 & \textbf{0.829} \\
Center    & 0.618          & 0.618 & \textbf{0.676} \\
Edge-Loc  & \textbf{0.705} & 0.682 & \textbf{0.705} \\
Edge-Ring & 0.941          & \textbf{0.971} & 0.961 \\
Donut     & 1.000          & 1.000 & 1.000 \\
Random    & 1.000          & 1.000 & 1.000 \\
Near-full & 1.000          & 1.000 & 1.000 \\
\midrule
\textbf{Overall recall} & ${\sim}$80.0\% & 82.4\% & \textbf{84.4\%} \\
\bottomrule
\end{tabular}
\end{table}

\textbf{Why 25\% masking outperforms 75\% on this domain.} The 75\% ratio is standard for natural images, which have dense, spatially redundant texture. Wafer maps are sparse binary grids with far less texture diversity. At 75\% masking, only 49 of 196 patches remain visible, providing insufficient local context for the model to learn meaningful wafer patch statistics. At 25\% masking, 147 patches remain visible per pass, giving the model richer local context and producing features better calibrated to wafer spatial structure.


\section{Ablation Studies and Variant Sweeps}
\label{app:ablations}

This appendix reports full sweep results for the ablation studies summarised in the main report.

% -------------------------------------------------------------------
\subsection{Phase 1: VAE KL-Weight Sweep}
\label{app:vae_sweep}

\begin{table}[H]
\centering
\small
\begin{tabularx}{\linewidth}{c X r r r}
\toprule
\textbf{KL weight} ($\beta$) & \textbf{Role} & \textbf{F1} & \textbf{AUROC} & \textbf{AUPRC} \\
\midrule
0.001 & Minimal regularisation  & \textbf{0.343} & 0.770 & 0.336 \\
0.005 & Balanced                & 0.340          & 0.772 & 0.372 \\
0.010 & Moderate regularisation & 0.335          & 0.766 & 0.369 \\
0.050 & Strong regularisation   & 0.302          & 0.750 & 0.329 \\
\bottomrule
\end{tabularx}
\caption{VAE KL-weight sweep. Increasing $\beta$ weakens reconstruction fidelity. All variants underperform the plain AE baseline (F1=0.468), confirming that the probabilistic objective does not help on this domain.}
\label{tab:vae_beta_sweep}
\end{table}

% -------------------------------------------------------------------
\subsection{Phase 1: Baseline AE Score Ablation}
\label{app:score_ablation}

\begin{figure}[H]
\centering
\includegraphics[width=0.95\linewidth]{images/autoencoder/ae64_score_ablation.png}
\caption{Score ablation on the baseline AE checkpoint: seven reduction rules applied to the same checkpoint produce F1 ranging from 0.24 to 0.47. Generated by \texttt{experiments/anomaly\_detection/autoencoder/x64/baseline/notebook.ipynb}.}
\label{fig:ae64_score_ablation}
\end{figure}

\begin{table}[H]
\centering
\small
\begin{tabularx}{\linewidth}{l X r}
\toprule
\textbf{Scoring Rule} & \textbf{Definition} & \textbf{F1} \\
\midrule
\texttt{topk\_abs\_mean} & Average top 20\% of pixel errors (sorted by magnitude) & \textbf{0.468} \\
\texttt{mse\_mean}       & Average squared pixel errors across entire image        & 0.410 \\
\texttt{max\_abs}        & Single worst (maximum absolute) pixel error             & 0.324 \\
\texttt{foreground\_mse} & Average error in foreground-masked region               & 0.317 \\
\texttt{mae\_mean}       & Average absolute pixel errors (MAE)                     & 0.300 \\
\texttt{pooled\_mae\_mean} & MAE with spatial pooling                              & 0.293 \\
\texttt{foreground\_mae} & MAE in foreground region                                & 0.239 \\
\bottomrule
\end{tabularx}
\caption{Score ablation: seven reduction rules on the same AE checkpoint. \texttt{topk\_abs\_mean} wins by $+0.057$ F1 over \texttt{mse\_mean}, confirming that selecting high-error pixels outperforms averaging all pixels.}
\label{tab:score_ablation}
\end{table}

\begin{figure}[H]
\centering
\includegraphics[width=0.95\linewidth]{images/autoencoder/score_ablation_summary.png}
\caption{Score ablation on x224 checkpoint: \texttt{topk\_abs\_mean} remains dominant at higher resolution, confirming that the pixel-selection insight holds across resolutions. Generated by \texttt{experiments/anomaly\_detection/autoencoder/x224/main/notebook.ipynb}.}
\label{fig:ae224_score_ablation}
\end{figure}

% -------------------------------------------------------------------
\subsection{Phase 3: Teacher-Student Backbone Comparison}
\label{app:ts_comparison}

\begin{table}[H]
\centering
\small
\caption{Teacher-student distillation results across all backbone variants.}
\label{tab:ts_results}
\begin{tabularx}{\linewidth}{>{\raggedright\arraybackslash}X c c c c}
\toprule
\textbf{Backbone} & \textbf{F1} & \textbf{AUROC} & \textbf{AUPRC} & \textbf{Notes} \\
\midrule
ResNet18 (topk20)              & 0.495          & 0.894          & 0.519          & Baseline \\
ResNet50 (mixed 2:1 topk20)    & \textbf{0.525} & 0.909          & 0.599          & Best CNN result \\
WideResNet50-2 layer2 (topk25) & 0.508          & 0.920          & 0.540          & Single-layer \\
WideResNet50-2 multilayer (topk15) & 0.524      & \textbf{0.923} & 0.546          & Multilayer \\
ViT-B/16 (x224)                & 0.163          & 0.661          & 0.106          & Architectural mismatch \\
\bottomrule
\end{tabularx}
\end{table}

\begin{figure}[H]
\centering
\includegraphics[width=\linewidth]{images/teacherstudent/ts_family_comparison_full.png}
\caption{Teacher-student distillation: all backbone variants compared across F1, AUROC, and AUPRC. WideResNet50-2 multilayer achieves the strongest AUROC (0.923); F1 is near-identical with ResNet50 (0.524 vs 0.525). Generated by \texttt{experiments/anomaly\_detection/report\_figures/notebook.ipynb}.}
\label{fig:ts_family_full}
\end{figure}

% -------------------------------------------------------------------
\subsection{Phase 3: FastFlow Architecture and Variants}
\label{app:fastflow}

\textbf{Stage 1: Feature extraction.} The input wafer map is converted from grayscale to RGB via channel repetition, resized to backbone input size via bilinear interpolation, then passed through a frozen WideResNet50-2 backbone. Features from layer2 (512 channels) and layer3 (2048 channels) are extracted and concatenated into a single tensor C.

\textbf{Stage 2: Normalising flow.} A flow with 4 sequential steps processes tensor C. Each step alternates between an orthogonal 1$\times$1 convolution (invertible channel mixing) and an affine coupling layer. The coupling layer splits channels in half, passes the first half through a small CNN (Conv3$\times$3 $\rightarrow$ GELU $\rightarrow$ Conv1$\times$1 $\rightarrow$ GELU $\rightarrow$ Conv3$\times$3, zero-initialised) to compute shift and log-scale parameters, then applies an affine transform to the second half.

\textbf{Stage 3: Anomaly scoring.} The transformed tensor Z is converted to a per-pixel anomaly map by computing $0.5 \times Z^2$ spatially. This heatmap is upsampled back to 64$\times$64 and reduced to a single wafer-level anomaly score via mean pooling.

\begin{figure}[H]
\centering
\includegraphics[width=\linewidth]{images/fastflow/fastflow_1.png}
\caption{FastFlow preprocessing and feature extraction pipeline. Generated by \texttt{experiments/anomaly\_detection/fastflow/x64/main/notebook.ipynb}.}
\label{fig:fastflow_pipeline}
\end{figure}

\begin{figure}[H]
\centering
\includegraphics[width=\linewidth]{images/fastflow/fastflow_2.png}
\caption{FastFlow normalising flow architecture: 4 sequential steps of orthogonal 1$\times$1 convolution and affine coupling. The final transformed tensor Z is converted to an anomaly map ($0.5 \times Z^2$). Generated by \texttt{experiments/anomaly\_detection/fastflow/x64/main/notebook.ipynb}.}
\label{fig:fastflow_flow}
\end{figure}

\begin{table}[H]
\centering
\small
\caption{FastFlow configuration variants on WideResNet50-2 backbone.}
\label{tab:fastflow_results}
\begin{tabularx}{\linewidth}{>{\raggedright\arraybackslash}X c c c c}
\toprule
\textbf{Configuration} & \textbf{F1} & \textbf{AUROC} & \textbf{AUPRC} & \textbf{Notes} \\
\midrule
Layer2+Layer3, 4 steps, mean      & \textbf{0.482} & \textbf{0.871} & \textbf{0.489} & Selected variant \\
Layer2+Layer3, 6 steps, mean      & 0.470          & 0.870          & 0.479          & More depth unhelpful \\
Layer2 only, 6 steps, topk-r0.15  & 0.442          & 0.884          & 0.460          & Single-layer weaker \\
\bottomrule
\end{tabularx}
\end{table}

\begin{figure}[H]
\centering
\includegraphics[width=\linewidth]{images/fastflow/variant_comparison_metrics.png}
\caption{FastFlow variant comparison: F1, AUROC, AUPRC across three configurations. Feature breadth (layer2+layer3) matters more than flow depth. Generated by \texttt{experiments/anomaly\_detection/fastflow/x64/main/notebook.ipynb}.}
\label{fig:fastflow_variants}
\end{figure}

% -------------------------------------------------------------------
\subsection{Phase 4: ViT-B/16 Block Depth Sweep}
\label{app:vit_block_sweep}

\begin{table}[H]
\centering
\caption{ViT-B/16 PatchCore block depth sweep at $224\!\times\!224$. All runs use identical
hyperparameters (5\% coreset, topk\_patch\_ratio 0.10, cosine distance). Block indices are
0-based over the 12-block encoder.}
\label{tab:vit_block_sweep}
\begin{tabular}{llcccc}
\toprule
Block & Depth    & F1             & AUROC          & AUPRC          & Recall \\
\midrule
3            & Early    & 0.550          & 0.934          & 0.622          & 0.732 \\
\textbf{6}   & \textbf{Mid} & \textbf{0.580} & \textbf{0.954} & 0.642     & \textbf{0.800} \\
9            & Mid-late & 0.580          & 0.943          & \textbf{0.677} & 0.788 \\
11           & Final    & 0.435          & 0.866          & 0.490          & 0.564 \\
\bottomrule
\end{tabular}
\end{table}

% -------------------------------------------------------------------
\subsection{Phase 4: DINOv2 ViT-B/14 Block Sweep}
\label{app:dinov2_block_sweep}

\begin{table}[H]
\centering
\caption{DINOv2 ViT-B/14 PatchCore block sweep at $224\!\times\!224$. All runs use identical hyperparameters.}
\label{tab:dinov2_block_sweep}
\begin{tabular}{llcccc}
\toprule
Block & Depth    & F1             & AUROC          & AUPRC          & Recall \\
\midrule
\textbf{6}   & \textbf{Mid} & \textbf{0.521} & \textbf{0.926} & 0.549     & \textbf{73.6\%} \\
9            & Mid-late & 0.492          & 0.915          & \textbf{0.561} & 68.4\% \\
11           & Final    & 0.443          & 0.885          & 0.504          & 58.0\% \\
\bottomrule
\end{tabular}
\end{table}

% -------------------------------------------------------------------
\subsection{Phase 4: MAE Fine-Tuning Setup and Per-Defect Results}
\label{app:mae_finetune}

\textbf{Setup.} The last 6 transformer blocks (blocks 6--11) and the final LayerNorm are fine-tuned for 10 epochs on the same 40k training normals used for the memory bank. A 2-layer MLP decoder (768 $\to$ 512 $\to$ patch pixels) produces the reconstruction target and is discarded after pre-training. AdamW with LR $=10^{-4}$ and cosine decay is used throughout. The 75\% masking run converges from loss 0.069 (epoch 1) to 0.042 (epoch 10); the 25\% masking run converges to 0.032 (epoch 10).

\begin{table}[H]
\centering
\caption{Frozen vs.\ MAE fine-tuned ViT-B/16 PatchCore at two masking ratios. All settings identical except backbone weights and masking ratio. All variants use the 95th-percentile validation-normal threshold.}
\label{tab:vit_mae_comparison}
\begin{tabular}{lcccc}
\toprule
Backbone & F1 & AUROC & AUPRC & Recall \\
\midrule
Frozen (ImageNet)          & 0.595          & 0.956          & 0.671          & ${\sim}$80.0\% \\
MAE fine-tuned (mask 75\%) & 0.595          & 0.959          & 0.717          & 82.4\% \\
MAE fine-tuned (mask 25\%) & \textbf{0.594} & \textbf{0.962} & \textbf{0.762} & \textbf{84.4\%} \\
\bottomrule
\end{tabular}
\end{table}

\begin{table}[H]
\centering
\caption{Per-defect recall for frozen ViT-B/16 vs.\ MAE fine-tuned variants.}
\label{tab:vit_mae_per_defect}
\begin{tabular}{lccc}
\toprule
\textbf{Defect type} & \textbf{Frozen} & \textbf{MAE 75\%} & \textbf{MAE 25\%} \\
\midrule
Scratch   & 0.727          & 0.455 & \textbf{0.636} \\
Loc       & \textbf{0.829} & 0.805 & \textbf{0.829} \\
Center    & 0.618          & 0.618 & \textbf{0.676} \\
Edge-Loc  & \textbf{0.705} & 0.682 & \textbf{0.705} \\
Edge-Ring & 0.941          & \textbf{0.971} & 0.961 \\
Donut     & 1.000          & 1.000 & 1.000 \\
Random    & 1.000          & 1.000 & 1.000 \\
Near-full & 1.000          & 1.000 & 1.000 \\
\midrule
\textbf{Overall recall} & ${\sim}$80.0\% & 82.4\% & \textbf{84.4\%} \\
\bottomrule
\end{tabular}
\end{table}

\textbf{Why 25\% masking outperforms 75\% on this domain.} The 75\% ratio is standard for natural images, which have dense, spatially redundant texture. Wafer maps are sparse binary grids with far less texture diversity. At 75\% masking, only 49 of 196 patches remain visible, providing insufficient local context for the model to learn meaningful wafer patch statistics. At 25\% masking, 147 patches remain visible per pass, giving the model richer local context and producing features better calibrated to wafer spatial structure.


\section{Ablation Studies and Variant Sweeps}
\label{app:ablations}

This appendix reports full sweep results for the ablation studies summarised in the main report.

% -------------------------------------------------------------------
\subsection{Phase 1: VAE KL-Weight Sweep}
\label{app:vae_sweep}

\begin{table}[H]
\centering
\small
\begin{tabularx}{\linewidth}{c X r r r}
\toprule
\textbf{KL weight} ($\beta$) & \textbf{Role} & \textbf{F1} & \textbf{AUROC} & \textbf{AUPRC} \\
\midrule
0.001 & Minimal regularisation  & \textbf{0.343} & 0.770 & 0.336 \\
0.005 & Balanced                & 0.340          & 0.772 & 0.372 \\
0.010 & Moderate regularisation & 0.335          & 0.766 & 0.369 \\
0.050 & Strong regularisation   & 0.302          & 0.750 & 0.329 \\
\bottomrule
\end{tabularx}
\caption{VAE KL-weight sweep. Increasing $\beta$ weakens reconstruction fidelity. All variants underperform the plain AE baseline (F1=0.468), confirming that the probabilistic objective does not help on this domain.}
\label{tab:vae_beta_sweep}
\end{table}

% -------------------------------------------------------------------
\subsection{Phase 1: Baseline AE Score Ablation}
\label{app:score_ablation}

\begin{figure}[H]
\centering
\includegraphics[width=0.95\linewidth]{images/autoencoder/ae64_score_ablation.png}
\caption{Score ablation on the baseline AE checkpoint: seven reduction rules applied to the same checkpoint produce F1 ranging from 0.24 to 0.47. Generated by \texttt{experiments/anomaly\_detection/autoencoder/x64/baseline/notebook.ipynb}.}
\label{fig:ae64_score_ablation}
\end{figure}

\begin{table}[H]
\centering
\small
\begin{tabularx}{\linewidth}{l X r}
\toprule
\textbf{Scoring Rule} & \textbf{Definition} & \textbf{F1} \\
\midrule
\texttt{topk\_abs\_mean} & Average top 20\% of pixel errors (sorted by magnitude) & \textbf{0.468} \\
\texttt{mse\_mean}       & Average squared pixel errors across entire image        & 0.410 \\
\texttt{max\_abs}        & Single worst (maximum absolute) pixel error             & 0.324 \\
\texttt{foreground\_mse} & Average error in foreground-masked region               & 0.317 \\
\texttt{mae\_mean}       & Average absolute pixel errors (MAE)                     & 0.300 \\
\texttt{pooled\_mae\_mean} & MAE with spatial pooling                              & 0.293 \\
\texttt{foreground\_mae} & MAE in foreground region                                & 0.239 \\
\bottomrule
\end{tabularx}
\caption{Score ablation: seven reduction rules on the same AE checkpoint. \texttt{topk\_abs\_mean} wins by $+0.057$ F1 over \texttt{mse\_mean}, confirming that selecting high-error pixels outperforms averaging all pixels.}
\label{tab:score_ablation}
\end{table}

\begin{figure}[H]
\centering
\includegraphics[width=0.95\linewidth]{images/autoencoder/score_ablation_summary.png}
\caption{Score ablation on x224 checkpoint: \texttt{topk\_abs\_mean} remains dominant at higher resolution, confirming that the pixel-selection insight holds across resolutions. Generated by \texttt{experiments/anomaly\_detection/autoencoder/x224/main/notebook.ipynb}.}
\label{fig:ae224_score_ablation}
\end{figure}

% -------------------------------------------------------------------
\subsection{Phase 3: Teacher-Student Backbone Comparison}
\label{app:ts_comparison}

\begin{table}[H]
\centering
\small
\caption{Teacher-student distillation results across all backbone variants.}
\label{tab:ts_results}
\begin{tabularx}{\linewidth}{>{\raggedright\arraybackslash}X c c c c}
\toprule
\textbf{Backbone} & \textbf{F1} & \textbf{AUROC} & \textbf{AUPRC} & \textbf{Notes} \\
\midrule
ResNet18 (topk20)              & 0.495          & 0.894          & 0.519          & Baseline \\
ResNet50 (mixed 2:1 topk20)    & \textbf{0.525} & 0.909          & 0.599          & Best CNN result \\
WideResNet50-2 layer2 (topk25) & 0.508          & 0.920          & 0.540          & Single-layer \\
WideResNet50-2 multilayer (topk15) & 0.524      & \textbf{0.923} & 0.546          & Multilayer \\
ViT-B/16 (x224)                & 0.163          & 0.661          & 0.106          & Architectural mismatch \\
\bottomrule
\end{tabularx}
\end{table}

\begin{figure}[H]
\centering
\includegraphics[width=\linewidth]{images/teacherstudent/ts_family_comparison_full.png}
\caption{Teacher-student distillation: all backbone variants compared across F1, AUROC, and AUPRC. WideResNet50-2 multilayer achieves the strongest AUROC (0.923); F1 is near-identical with ResNet50 (0.524 vs 0.525). Generated by \texttt{experiments/anomaly\_detection/report\_figures/notebook.ipynb}.}
\label{fig:ts_family_full}
\end{figure}

% -------------------------------------------------------------------
\subsection{Phase 3: FastFlow Architecture and Variants}
\label{app:fastflow}

\textbf{Stage 1: Feature extraction.} The input wafer map is converted from grayscale to RGB via channel repetition, resized to backbone input size via bilinear interpolation, then passed through a frozen WideResNet50-2 backbone. Features from layer2 (512 channels) and layer3 (2048 channels) are extracted and concatenated into a single tensor C.

\textbf{Stage 2: Normalising flow.} A flow with 4 sequential steps processes tensor C. Each step alternates between an orthogonal 1$\times$1 convolution (invertible channel mixing) and an affine coupling layer. The coupling layer splits channels in half, passes the first half through a small CNN (Conv3$\times$3 $\rightarrow$ GELU $\rightarrow$ Conv1$\times$1 $\rightarrow$ GELU $\rightarrow$ Conv3$\times$3, zero-initialised) to compute shift and log-scale parameters, then applies an affine transform to the second half.

\textbf{Stage 3: Anomaly scoring.} The transformed tensor Z is converted to a per-pixel anomaly map by computing $0.5 \times Z^2$ spatially. This heatmap is upsampled back to 64$\times$64 and reduced to a single wafer-level anomaly score via mean pooling.

\begin{figure}[H]
\centering
\includegraphics[width=\linewidth]{images/fastflow/fastflow_1.png}
\caption{FastFlow preprocessing and feature extraction pipeline. Generated by \texttt{experiments/anomaly\_detection/fastflow/x64/main/notebook.ipynb}.}
\label{fig:fastflow_pipeline}
\end{figure}

\begin{figure}[H]
\centering
\includegraphics[width=\linewidth]{images/fastflow/fastflow_2.png}
\caption{FastFlow normalising flow architecture: 4 sequential steps of orthogonal 1$\times$1 convolution and affine coupling. The final transformed tensor Z is converted to an anomaly map ($0.5 \times Z^2$). Generated by \texttt{experiments/anomaly\_detection/fastflow/x64/main/notebook.ipynb}.}
\label{fig:fastflow_flow}
\end{figure}

\begin{table}[H]
\centering
\small
\caption{FastFlow configuration variants on WideResNet50-2 backbone.}
\label{tab:fastflow_results}
\begin{tabularx}{\linewidth}{>{\raggedright\arraybackslash}X c c c c}
\toprule
\textbf{Configuration} & \textbf{F1} & \textbf{AUROC} & \textbf{AUPRC} & \textbf{Notes} \\
\midrule
Layer2+Layer3, 4 steps, mean      & \textbf{0.482} & \textbf{0.871} & \textbf{0.489} & Selected variant \\
Layer2+Layer3, 6 steps, mean      & 0.470          & 0.870          & 0.479          & More depth unhelpful \\
Layer2 only, 6 steps, topk-r0.15  & 0.442          & 0.884          & 0.460          & Single-layer weaker \\
\bottomrule
\end{tabularx}
\end{table}

\begin{figure}[H]
\centering
\includegraphics[width=\linewidth]{images/fastflow/variant_comparison_metrics.png}
\caption{FastFlow variant comparison: F1, AUROC, AUPRC across three configurations. Feature breadth (layer2+layer3) matters more than flow depth. Generated by \texttt{experiments/anomaly\_detection/fastflow/x64/main/notebook.ipynb}.}
\label{fig:fastflow_variants}
\end{figure}

% -------------------------------------------------------------------
\subsection{Phase 4: ViT-B/16 Block Depth Sweep}
\label{app:vit_block_sweep}

\begin{table}[H]
\centering
\caption{ViT-B/16 PatchCore block depth sweep at $224\!\times\!224$. All runs use identical
hyperparameters (5\% coreset, topk\_patch\_ratio 0.10, cosine distance). Block indices are
0-based over the 12-block encoder.}
\label{tab:vit_block_sweep}
\begin{tabular}{llcccc}
\toprule
Block & Depth    & F1             & AUROC          & AUPRC          & Recall \\
\midrule
3            & Early    & 0.550          & 0.934          & 0.622          & 0.732 \\
\textbf{6}   & \textbf{Mid} & \textbf{0.580} & \textbf{0.954} & 0.642     & \textbf{0.800} \\
9            & Mid-late & 0.580          & 0.943          & \textbf{0.677} & 0.788 \\
11           & Final    & 0.435          & 0.866          & 0.490          & 0.564 \\
\bottomrule
\end{tabular}
\end{table}

% -------------------------------------------------------------------
\subsection{Phase 4: DINOv2 ViT-B/14 Block Sweep}
\label{app:dinov2_block_sweep}

\begin{table}[H]
\centering
\caption{DINOv2 ViT-B/14 PatchCore block sweep at $224\!\times\!224$. All runs use identical hyperparameters.}
\label{tab:dinov2_block_sweep}
\begin{tabular}{llcccc}
\toprule
Block & Depth    & F1             & AUROC          & AUPRC          & Recall \\
\midrule
\textbf{6}   & \textbf{Mid} & \textbf{0.521} & \textbf{0.926} & 0.549     & \textbf{73.6\%} \\
9            & Mid-late & 0.492          & 0.915          & \textbf{0.561} & 68.4\% \\
11           & Final    & 0.443          & 0.885          & 0.504          & 58.0\% \\
\bottomrule
\end{tabular}
\end{table}

% -------------------------------------------------------------------
\subsection{Phase 4: MAE Fine-Tuning Setup and Per-Defect Results}
\label{app:mae_finetune}

\textbf{Setup.} The last 6 transformer blocks (blocks 6--11) and the final LayerNorm are fine-tuned for 10 epochs on the same 40k training normals used for the memory bank. A 2-layer MLP decoder (768 $\to$ 512 $\to$ patch pixels) produces the reconstruction target and is discarded after pre-training. AdamW with LR $=10^{-4}$ and cosine decay is used throughout. The 75\% masking run converges from loss 0.069 (epoch 1) to 0.042 (epoch 10); the 25\% masking run converges to 0.032 (epoch 10).

\begin{table}[H]
\centering
\caption{Frozen vs.\ MAE fine-tuned ViT-B/16 PatchCore at two masking ratios. All settings identical except backbone weights and masking ratio. All variants use the 95th-percentile validation-normal threshold.}
\label{tab:vit_mae_comparison}
\begin{tabular}{lcccc}
\toprule
Backbone & F1 & AUROC & AUPRC & Recall \\
\midrule
Frozen (ImageNet)          & 0.595          & 0.956          & 0.671          & ${\sim}$80.0\% \\
MAE fine-tuned (mask 75\%) & 0.595          & 0.959          & 0.717          & 82.4\% \\
MAE fine-tuned (mask 25\%) & \textbf{0.594} & \textbf{0.962} & \textbf{0.762} & \textbf{84.4\%} \\
\bottomrule
\end{tabular}
\end{table}

\begin{table}[H]
\centering
\caption{Per-defect recall for frozen ViT-B/16 vs.\ MAE fine-tuned variants.}
\label{tab:vit_mae_per_defect}
\begin{tabular}{lccc}
\toprule
\textbf{Defect type} & \textbf{Frozen} & \textbf{MAE 75\%} & \textbf{MAE 25\%} \\
\midrule
Scratch   & 0.727          & 0.455 & \textbf{0.636} \\
Loc       & \textbf{0.829} & 0.805 & \textbf{0.829} \\
Center    & 0.618          & 0.618 & \textbf{0.676} \\
Edge-Loc  & \textbf{0.705} & 0.682 & \textbf{0.705} \\
Edge-Ring & 0.941          & \textbf{0.971} & 0.961 \\
Donut     & 1.000          & 1.000 & 1.000 \\
Random    & 1.000          & 1.000 & 1.000 \\
Near-full & 1.000          & 1.000 & 1.000 \\
\midrule
\textbf{Overall recall} & ${\sim}$80.0\% & 82.4\% & \textbf{84.4\%} \\
\bottomrule
\end{tabular}
\end{table}

\textbf{Why 25\% masking outperforms 75\% on this domain.} The 75\% ratio is standard for natural images, which have dense, spatially redundant texture. Wafer maps are sparse binary grids with far less texture diversity. At 75\% masking, only 49 of 196 patches remain visible, providing insufficient local context for the model to learn meaningful wafer patch statistics. At 25\% masking, 147 patches remain visible per pass, giving the model richer local context and producing features better calibrated to wafer spatial structure.
